\section{Optimality for a fixed defect kernel}

The vacuum-derived weight is canonical when a physical \(T\) is available.
For a fixed bad kernel \(Q\), a separate variational statement identifies the
best possible positive diagonal weight.

\begin{theorem}[Collatz--Wielandt defect optimum]\label{thm:cw}
Let \(Q\) be a finite irreducible nonnegative matrix.  Then
\begin{equation}\label{eq:cw}
 \rhoPF(Q)
 =\inf_{V>0}\max_i\frac{(QV)_i}{V_i}
 =\sup_{V>0}\min_i\frac{(QV)_i}{V_i}.
\end{equation}
The infimum is achieved by a positive Perron vector
\(V_{\mathrm{PF}}\), for which
\(QV_{\mathrm{PF}}=\rhoPF(Q)V_{\mathrm{PF}}\).
\end{theorem}

Thus no positive weighting can produce a ceiling below the spectral radius of
the declared bad subkernel.  This is an exact obstruction as well as an
optimization theorem.

\begin{warning}
Computing \(V_{\mathrm{PF}}\) from a finite kernel does not prove a uniform
Yang--Mills estimate.  The kernel must first be the complete normalized
physical conditional kernel.  Its state space, positivity, boundary
dependence, topology, and cutoff dependence cannot be inferred from Perron
theory.
\end{warning}

For an exploration or RG kernel, \(V_{\mathrm{PF}}\) is a Lyapunov
eigenweight.  Calling it a vacuum wavefunction is justified only after an
independent theorem identifies that kernel with reflection-positive temporal
transfer.
